\chapter{Integrali di linea di $1^\circ$ specie}

\section{Arco di curva regolare}
Le curve sono definite da funzioni a valori vettoriali. Un \textbf{arco di curva} si dice \textbf{regolare} se ammette una parametrizzazione $\mathbf{x}(t)$ con $t \in [a,b]$ che soddisfa le seguenti condizioni:
\begin{itemize}
    \item è iniettiva;
    \item è derivabile (con derivata continua e mai nulla).
\end{itemize}

La \textbf{lunghezza dell'arco di curva} è data dalla formula:
\[
l = \int_{a}^{b} \sqrt{\left(\frac{dx_1}{dt}\right)^2 + \dots + \left(\frac{dx_n}{dt}\right)^2} \, dt = \int_{a}^{b} \left\| \frac{d\vec{r}}{dt} \right\| \, dt
\]

\section{Ascissa curvilinea}
Per parametrizzare una curva ci sono infiniti modi, ossia si possono scegliere infiniti parametri arbitrari che soddisfano le condizioni di regolarità. Di fatto, esiste un parametro privilegiato chiamato \textbf{ascissa curvilinea} (o parametro naturale) $s$.

Data una curva parametrizzata, fissando un punto iniziale su di essa $\vec{r}(t_0)=P_0$, l'ascissa curvilinea rappresenta la lunghezza percorsa lungo la curva a partire da $P_0$:
\[
l(P_0) = \int_{t_0}^{t} \left\| \frac{d\vec{r}}{dt} \right\| \, dt
\]
Possiamo quindi sostituire $\vec{r}(t)$ e riparametrizzare la curva in funzione dell'ascissa curvilinea, ottenendo $\vec{l}$.

\section{Definizione di Integrale di linea di $1^\circ$ specie}
Sia $\gamma$ una curva regolare (definita in un intervallo $[a,b]$) e sia $f: \mathbb{R}^n \rightarrow \mathbb{R}$ una funzione continua. 
L'integrale di linea di $1^\circ$ specie valuta la restrizione della funzione $f$ lungo la curva $\gamma$.

Considerando la funzione composta lungo la curva $F(l) = f(x_1(l), \dots, x_n(l))$, l'integrale si definisce formalmente rispetto all'ascissa curvilinea $s$ (con lunghezza totale $L$) come:
\[
\int_{\gamma} f \, ds = \int_{0}^{L} f(\mathbf{x}(s)) \, ds
\]
Nella pratica, passando a una parametrizzazione generica $t \in [a,b]$ e ricordando che 
\[
dl = \frac{d\vec{r}}{dt} dt
\]
la formula per il calcolo diventa:
\[
\int_{\gamma} f \, dl = \int_a^b F(t) \, \frac{dl}{dt}dt = \int_a^b F(l(t)) \, ||\frac{d\vec{r}}{dt}dt|| = \int_a^b F(t) \, dt
\]

\section{Esercizio: Calcola l'integrale di linea di 1ª specie} 
\[
\int_{\gamma} (x+y) \, ds
\]
dove $\gamma$ è il segmento di retta che congiunge l'origine $(0,0)$ al punto $(1,1)$.

\vspace{0.5cm}

Parametrizziamo il segmento di retta per $t \in [0,1]$:
\[
\vec{r}(t) = 
\begin{cases} 
x(t) = 0 + t(1-0) = t \\ 
y(t) = 0 + t(1-0) = t 
\end{cases}
\]

Calcoliamo il vettore derivato:
\[
\frac{d\vec{r}}{dt} = \vec{r}'(t) = \begin{pmatrix} 1 \\ 1 \end{pmatrix}
\]

Troviamo il differenziale d'arco $ds$:
\[
ds = \|\vec{r}'(t)\| \, dt = \sqrt{1^2 + 1^2} \, dt = \sqrt{2} \, dt
\]

Sostituiamo nell'integrale e risolviamo:
\begin{align*}
\int_{\gamma} (x+y) \, ds &= \int_{0}^{1} (t + t) \sqrt{2} \, dt \\[1ex]
&= \int_{0}^{1} 2t\sqrt{2} \, dt \\[1ex]
&= 2\sqrt{2} \int_{0}^{1} t \, dt \\[1ex]
&= 2\sqrt{2} \left[ \frac{t^2}{2} \right]_{0}^{1} \\[1ex]
&= 2\sqrt{2} \cdot \frac{1}{2} \\[1ex]
&= \sqrt{2}
\end{align*}